
    $x$ is an interior point of $X$, that is, there exists
$\varepsilon > 0$ such that
\[
\|z-x\| \leq \varepsilon \Longrightarrow z \in X.
\]
It means that the (convex) ball of radius $\varepsilon$ and centered
at $x$ is completely in $X$.

Consider an arbitrary direction $d\neq 0$, and the point
\[
z=x+\varepsilon \frac{d}{\parallel d \parallel}.
\]
The point $z$ is at the intersection of the direction $d$ and the border
of the ball.
As
\[
\|z-x\| = \left\|\varepsilon \frac{d}{\parallel d \parallel}\right\| = \varepsilon,
\]
the definition of the interior point guarantees that  $z \in X$. By
convexity of the ball, all points $x + \alpha d$ with
\[
0 \leq \alpha \leq \frac{\varepsilon}{\|d\|}
\]
also belong to the ball, and therefore, to $X$, proving that the
direction $d$ is feasible.

\begin{center}
\begin{tikzpicture}[scale=2]
  \draw (0,0) ellipse[x radius=3,y radius= 1.5];
  \node[at={(1,1.6)}] {$X$} ;
  \node[at={(1,0)},circle,fill,label=$x$] {} ;
  \node[at={(2,0)},circle,fill,label=$z$] {} ;
  \draw (1,0) circle[radius=1];
\draw [Circle-] (1,0) -- node[left,red] {$\varepsilon$} (60:1);  
  \draw[-latex] (1,0) -- (2.3,0)  node[above] {$d$};
  \draw[-latex] (1,0) -- (2.3,0)  node[above] {$d$};
\end{tikzpicture}
\end{center}
